\documentclass[preview]{standalone}

\usepackage{amsmath}
\usepackage{amssymb}
\usepackage{stellar}
\usepackage{definitions}
\usepackage{tikz}
\usetikzlibrary{cd}
\usepackage{bettelini}

\begin{document}

\id{series-of-functions-convergence-criteria}
\genpage

\section{Total convergence}

\begin{snippetdefinition}{function-series-total-convergence-definition}{Total convergence}
    Let \(S\subseteq\realnumbers\) and let
    \(f_n\colon S\fromto\realnumbers\) be bounded for every
    \(n\in\naturalnumbers\). The \series \(\sum_{n=1}^{\infty}f_n\)
    \emph{converges totally} on \(S\) if the numerical series
    \[
        \sum_{n=1}^{\infty}\lVert f_n\rVert_{\infty,S},
        \qquad
        \lVert f_n\rVert_{\infty,S}=\sup_{x\in S}|f_n(x)|,
    \]
    converges.
\end{snippetdefinition}

\begin{snippet}{function-series-total-convergence-majorant-characterization}
    A \series \(\sum f_n\) converges totally on \(S\) if and only if there
    exists a summable numerical sequence \(\{M_n\}\) such that
    \[
        \lVert f_n\rVert_{\infty,S}\leq M_n,
        \qquad \forall n\in\naturalnumbers.
    \]
    Indeed, one implication follows by comparison, while for the converse one
    can choose \(M_n=\lVert f_n\rVert_{\infty,S}\).
\end{snippet}

\begin{snippettheorem}{weierstrass-function-series-test}{Weierstrass test}
    Let \(S\subseteq\realnumbers\) and let
    \(f_n\colon S\fromto\realnumbers\). Suppose there exists a sequence of
    nonnegative real numbers \(\{M_n\}_{n\in\naturalnumbers}\) such that
    \[
        |f_n(x)|\leq M_n,
        \qquad \forall x\in S,\ \forall n\in\naturalnumbers,
    \]
    and \(\sum_{n=1}^{\infty}M_n<\infty\). Then
    \(\sum_{n=1}^{\infty}f_n\) converges uniformly on \(S\).
\end{snippettheorem}

\begin{snippetproof}{weierstrass-function-series-test-proof}{weierstrass-function-series-test}{Weierstrass test}
    Given \(\varepsilon>0\), the numerical Cauchy criterion supplies
    \(N\in\naturalnumbers\) such that
    \[
        \sum_{k=n+1}^{m}M_k<\varepsilon,
        \qquad \forall m>n\geq N.
    \]
    Therefore, for every \(x\in S\),
    \[
        \left|\sum_{k=n+1}^{m}f_k(x)\right|
        \leq\sum_{k=n+1}^{m}|f_k(x)|
        \leq\sum_{k=n+1}^{m}M_k<\varepsilon.
    \]
    The uniform Cauchy criterion completes the proof.
\end{snippetproof}

\begin{snippetcorollary}{function-series-total-implies-uniform-convergence}{Total convergence implies uniform convergence}
    Let \(S\subseteq\realnumbers\) and let
    \(f_n\colon S\fromto\realnumbers\). If
    \(\sum_{n=1}^{\infty}f_n\) converges totally on \(S\), then it converges
    uniformly on \(S\).
\end{snippetcorollary}

\begin{snippetproof}{function-series-total-implies-uniform-convergence-proof}{function-series-total-implies-uniform-convergence}{Total convergence implies uniform convergence}
    Total convergence means that
    \(\sum_n\lVert f_n\rVert_{\infty,S}<\infty\). Thus, for every
    \(\varepsilon>0\), there exists \(N\) such that, whenever \(m>n\geq N\),
    \[
        \left|\sum_{k=n+1}^{m}f_k(x)\right|
        \leq\sum_{k=n+1}^{m}|f_k(x)|
        \leq\sum_{k=n+1}^{m}\lVert f_k\rVert_{\infty,S}
        <\varepsilon
    \]
    for every \(x\in S\). The uniform Cauchy criterion applies.
\end{snippetproof}

\section{Comparison of convergence notions}

\begin{snippet}{function-series-convergence-implications}
    The principal implications between convergence notions are
    \[
        \begin{tikzcd}[column sep=large]
            \text{Total convergence}
                \arrow[r, Rightarrow]
                & \text{Uniform absolute convergence}
                \arrow[r, Rightarrow]
                \arrow[d, Rightarrow]
                & \text{Pointwise absolute convergence}
                \arrow[d, Rightarrow] \\
            & \text{Uniform convergence}
                \arrow[r, Rightarrow]
                & \text{Pointwise convergence.}
        \end{tikzcd}
    \]
\end{snippet}

\begin{snippetexample}{function-series-uniform-without-total-convergence-example}{Uniform convergence does not imply total convergence}
    On \(S=[1,+\infty)\), define
    \[
        f_n(x)=\frac1n\chi_{[n,n+1)}(x).
    \]
    Since the supports are pairwise disjoint, for \(m>n\),
    \[
        \sup_{x\in S}\left|\sum_{k=n+1}^{m}f_k(x)\right|
        =\max_{n<k\leq m}\frac1k
        =\frac1{n+1}\longrightarrow0.
    \]
    Hence the series converges uniformly. Its pointwise sum is the step
    function
    \[
        f(x)=\frac1{\lfloor x\rfloor},
        \qquad x\in[1,+\infty).
    \]
    However,
    \[
        \sum_{n=1}^{\infty}\lVert f_n\rVert_{\infty,S}
        =\sum_{n=1}^{\infty}\frac1n=+\infty,
    \]
    so the convergence is not total.
\end{snippetexample}

\begin{snippetproposition}{uniformly-convergent-function-series-term-test}{Term test for uniformly convergent series of functions}
    Let \(S\subseteq\realnumbers\) and let
    \(f_n\colon S\fromto\realnumbers\). If the \series
    \(\sum_{n=1}^{\infty}f_n\) converges uniformly on \(S\), then
    \(f_n\to0\) uniformly on \(S\).
\end{snippetproposition}

\begin{snippetproof}{uniformly-convergent-function-series-term-test-proof}{uniformly-convergent-function-series-term-test}{Term test for uniformly convergent series of functions}
    Let \(S_n=\sum_{k=1}^{n}f_k\). Uniform convergence of the series means
    that \(S_n\to S\) uniformly. Since \(f_n=S_n-S_{n-1}\),
    \[
        \lVert f_n\rVert_{\infty,S}
        \leq\lVert S_n-S\rVert_{\infty,S}
        +\lVert S_{n-1}-S\rVert_{\infty,S}\longrightarrow0.
    \]
\end{snippetproof}

\begin{snippetexercise}{uniform-term-convergence-not-sufficient-exercise}{Uniform convergence of the terms is not sufficient}
    On a nonempty set \(S\), let \(f_n(x)=1/n\) for every \(x\in S\).
    Verify that \(f_n\to0\) uniformly while
    \(\sum_{n=1}^{\infty}f_n\) does not converge pointwise, and therefore
    cannot converge uniformly.
\end{snippetexercise}

\begin{snippetexercise}{zeta-function-series-not-uniform-exercise}{A pointwise but not uniformly convergent series}
    Consider
    \[
        \sum_{n=1}^{\infty}\frac1{n^x},
        \qquad x\in(1,+\infty).
    \]
    Show that it converges pointwise and that its terms tend uniformly to zero,
    since
    \[
        \sup_{x>1}\frac1{n^x}=\frac1n\longrightarrow0,
    \]
    but the series does not converge uniformly. Indeed, uniform convergence
    would allow the limit \(x\to1^+\) to be interchanged with the limit of the
    partial sums, forcing the harmonic partial sums
    \(\sum_{k=1}^{n}1/k\) to converge.
\end{snippetexercise}

\begin{snippetproposition}{boundary-divergence-obstructs-uniform-series-convergence}{Boundary divergence obstructs uniform convergence}
    Let \(a<b\) and let \(f_n\colon[a,b)\fromto\realnumbers\) be continuous.
    If \(\sum_n f_n(x)\) converges pointwise for every \(x\in(a,b)\), while
    the numerical series \(\sum_n f_n(a)\) diverges, then
    \(\sum_n f_n\) does not converge uniformly on \((a,b)\).
\end{snippetproposition}

\begin{snippetproof}{boundary-divergence-obstructs-uniform-series-convergence-proof}{boundary-divergence-obstructs-uniform-series-convergence}{Boundary divergence obstructs uniform convergence}
    If the series converged uniformly on \((a,b)\), its sequence of partial
    sums \(S_n\) would converge uniformly there. Every \(S_n\) has the finite
    limit
    \[
        \lim_{x\to a^+}S_n(x)=\sum_{k=1}^{n}f_k(a).
    \]
    The interchange-of-limits theorem would imply that these boundary values
    form a convergent numerical sequence, contradicting divergence of
    \(\sum_n f_n(a)\).
\end{snippetproof}

\section{Alternating series}

\begin{snippetproposition}{uniform-alternating-function-series-test}{Uniform alternating-series test}
    Let \(S\subseteq\realnumbers\) and let
    \(f_n\colon S\fromto\realnumbers^{\geq0}\). Suppose
    \[
        f_{n+1}(x)\leq f_n(x),
        \qquad \forall x\in S,\ \forall n\in\naturalnumbers,
    \]
    and that \(f_n\to0\) uniformly on \(S\). Then the alternating series
    \[
        \sum_{n=1}^{\infty}(-1)^nf_n(x)
    \]
    converges uniformly on \(S\).
\end{snippetproposition}

\begin{snippetproof}{uniform-alternating-function-series-test-proof}{uniform-alternating-function-series-test}{Uniform alternating-series test}
    The remainder estimate for alternating numerical series, applied at each
    \(x\in S\), gives
    \[
        \left|
            \sum_{k=1}^{\infty}(-1)^kf_k(x)
            -\sum_{k=1}^{n}(-1)^kf_k(x)
        \right|\leq f_{n+1}(x).
    \]
    Taking the supremum over \(S\) and using uniform convergence
    \(f_{n+1}\to0\) proves uniform convergence of the series.
\end{snippetproof}

\begin{snippetexample}{function-series-uniform-without-absolute-convergence-example}{Uniform convergence does not imply absolute convergence}
    On any nonempty set \(S\), let \(f_n(x)=(-1)^n/n\). This is a series of
    constant functions, so convergence of the alternating harmonic series is
    uniform on \(S\). Its series of absolute values is the harmonic series and
    diverges at every point.
\end{snippetexample}

\begin{snippetexample}{function-series-absolute-uniform-without-total-convergence-example}{Absolute and uniform convergence do not imply total convergence}
    On \(S=\realnumbers^{>0}\), define
    \[
        f_n(x)=\frac1n\chi_{[n-1,n)}(x).
    \]
    At each point at most one term is nonzero, so the series converges
    absolutely pointwise. Moreover, for \(m>n\),
    \[
        \sup_{x\in S}\left|\sum_{k=n+1}^{m}f_k(x)\right|
        =\frac1{n+1}\longrightarrow0,
    \]
    hence it converges uniformly. Nevertheless,
    \[
        \sum_{n=1}^{\infty}\lVert f_n\rVert_{\infty,S}
        =\sum_{n=1}^{\infty}\frac1n=+\infty,
    \]
    so it does not converge totally.
\end{snippetexample}

\end{document}
